% 318_Masseableitung_Geltungsbereich_De_ch.tex
% Geltungsbereich der Masseableitungen in der FFGFT:
% Ruhemassen, Hadrondynamik und offene Brücken
% Dok. 318, August 2026


\section*{Zeichenerklärung}

\begin{center}
\begin{tabular}{lll}
\toprule
Symbol & Bedeutung & Wert / Einheit \\
\midrule
$M_Z$ & Masse des $Z$-Bosons; Referenzskala für $\alpha_s$ & $91{,}19$~GeV \\
$\Lambda_{\mathrm{QCD}}$ & QCD-Konfinementskala & $\approx 200$--$217$~MeV \\
$\alpha_s(\mu)$ & Starke Kopplung bei Skala $\mu$ & $\alpha_s(M_Z) \approx 0{,}118$; $\alpha_s(m_\tau) \approx 0{,}33$ \\
$D = N_c$ & Anzahl der Farbfreiheitsgrade (Dok.~160) & $D = 3$ \\
$y_e$ & Yukawa-Kopplung des Elektrons & $m_e/v \approx 2{,}08 \cdot 10^{-6}$ \\
$v$ & Higgs-VEV & $246$~GeV \\
$m_p, m_n$ & Proton-/Neutronmasse & $938{,}3$ / $939{,}6$~MeV \\
RGE & Renormierungsgruppengleichung & beschreibt $\mu$-Abhängigkeit von $\alpha_s$ \\
$K_{\text{frak}}$ & Fraktaler Interface-Faktor (eingerollt $\to$ SI) & $1-100\xi \approx 0{,}987$ \\
$m^{\text{bare}}$ / $m^{\text{SI}}$ & Geometrische / gemessene Masse & nat.~Einh. / MeV \\
$\xi$ & FFGFT-Geometrieparameter & $4/30000 \approx 1{,}333 \cdot 10^{-4}$ \\
\bottomrule
\end{tabular}
\end{center}

\section*{Zusammenfassung}

Die FFGFT leitet Teilchenmassen aus dem Geometrieparameter
$\xi = 4/30000$ auf dem kompakten 4D-Torus $T^4/\mathbb{Z}_3$ ab.
Dieser Geltungsbereich ist präzise für Ruhemassen aus dem
Higgs-Yukawa-Sektor (Leptonen, Quark-Ruhemassen). Hadronmassen wie die
Protonmasse enthalten einen dominant nichtperturbativen QCD-Beitrag,
der einen eigenen Behandlungsrahmen erfordert. Dieses Dokument zieht
die Grenze explizit und begründet sie physikalisch aus der
relativistischen Energierelation $E^2 = (m_0 c^2)^2 + (pc)^2$ und dem
Bild masseloser Gluonen im Protoninneren.

\begin{keyresult}[Statusmarkierung]
\textbf{[K]} FFGFT leitet Ruhemassen (Leptonen, Quark-Yukawa) aus
$\xi$-Geometrie ab.\quad \textbf{Hadronmassen} haben zusätzlich einen
nichtperturbativen QCD-Beitrag: \textbf{offene Brücke}.
\end{keyresult}

\section{Ausgangspunkt: die vollständige Energierelation}

Das populäre $E = m_0 c^2$ ist eine Näherung. Sie folgt aus der
Linearisierung des Lorentz-Faktors
\begin{equation}
  \gamma = \left(1 - \frac{v^2}{c^2}\right)^{-1/2}
  \approx 1 + \frac{1}{2}\frac{v^2}{c^2}
  \qquad (v \ll c)
\end{equation}
und gilt nur für kleine Geschwindigkeiten. Die vollständige,
exakte Relation lautet:
\begin{equation}
  \boxed{E^2 = (m_0 c^2)^2 + (pc)^2}
  \label{eq:volle_energie}
\end{equation}

Diese Gleichung ist der eigentliche Kern -- nicht die vereinfachte
Formel. Der zweite Term $pc$ verschwindet nicht, wenn ein Teilchen
masselose ist. Für masselose Teilchen ($m_0 = 0$) gilt
\begin{equation}
  E = pc,
\end{equation}
also trägt ein masseloses Teilchen, das sich mit Lichtgeschwindigkeit
bewegt, trotzdem zur Gesamtenergie und damit zur messbaren Masse
eines Systems bei.

\section{Das Proton: Masse als emergente Eigenschaft}

Die Protonmasse $m_p \approx 938{,}3$~MeV ist das klarste Beispiel
dieser Emergenz:

\begin{table}[h]
\centering
\caption{Energiebeiträge zur Protonmasse}
\label{tab:proton}
\begin{tabular}{llll}
\toprule
Beitrag & Wert & Anteil & Mechanismus \\
\midrule
$u$- und $d$-Quark-Ruhemassen & $\approx 10$~MeV & $\approx 1\,\%$
  & Higgs-Yukawa-Kopplung \\
Gluonen-Bewegungsenergie & $\approx 928$~MeV & $\approx 99\,\%$
  & QCD-Konfinement \\
\midrule
Protonmasse & $938{,}3$~MeV & $100\,\%$ & \\
\bottomrule
\end{tabular}
\end{table}

Das Proton besteht aus drei Quarks ($uud$), deren Ruhemassen zusammen
etwa $10$~MeV betragen. Der Rest -- $928$~MeV, rund $99\,\%$ -- kommt
von Gluonen, den Trägern der starken Kraft. Gluonen sind masselos
($m_0 = 0$). Sie bewegen sich mit Lichtgeschwindigkeit. Nach
Gleichung~(\ref{eq:volle_energie}) tragen sie trotzdem über ihren
Impuls zur Gesamtenergie des Systems bei.

\subsection{Zeitlosigkeit im Protoninneren}

Für masselose Teilchen, die sich mit Lichtgeschwindigkeit bewegen,
gilt maximale Zeitdilatation: Die Eigenzeit dieser Teilchen ist null.
Gluonen im Protoninneren existieren in einem Zustand, für den keine
Zeit vergeht. Die Ausdehnung der Welt ist für sie null --
Längenkontraktion ist maximal.

Und dennoch: Diese zeitlosen, masselosen Bestandteile bedingen den
Hauptteil dessen, was außen als Protonmasse gemessen wird. Masse ist
eine Eigenschaft an der Oberfläche eines Systems, nicht eine Grundeigenschaft
seiner Bestandteile.

\textbf{[S]} Diese Beobachtung berührt die Fundamentalrelation der
FFGFT direkt. $\tilde{T} \cdot m = 1$ besagt: Intrinsische Zeit und
Masse sind dual. Im Protoninneren liegt genau die zeitlose Ebene vor --
und diese bildet den Großteil der nach außen erscheinenden Masse.
Ruhemasse entsteht durch Projektion aus dem Zeitlosen ins Zeitbehaftete.
Das Proton illustriert $\tilde{T} \cdot m = 1$ unbeabsichtigt präzise.
Diese Verbindung ist konzeptuell belastbar, aber noch keine formale
Ableitung; sie ist für ein eigenes Dokument vorgemerkt.

\section{Zwei Sektoren -- zwei Mechanismen}

Das Verhältnis $m_p/m_e \approx 1836$ verbindet zwei physikalisch
verschiedene Sektoren:

\begin{center}
\begin{tabular}{lll}
\toprule
Größe & Sektor & Skala \\
\midrule
$m_e, m_\mu, m_\tau$ & Higgs-Yukawa & $\xi^{p_i} \cdot r_i \cdot v$ \\
$m_{\nu_e}, m_{\nu_\mu}, m_{\nu_\tau}$ & Higgs-Yukawa (supprimiert) & $\xi^{q_i} \cdot v$ \\
$m_u, m_d, m_s$ (Ruhe) & Higgs-Yukawa & (offen, Dok.~006) \\
$m_p, m_n$ (gesamt) & QCD-Konfinement + Higgs-Yukawa
  & $\Lambda_{\mathrm{QCD}} + m_{\mathrm{Quarks}}$ \\
\bottomrule
\end{tabular}
\end{center}

\textbf{[B]} Protonmasse und Elektronmasse entstehen durch völlig
verschiedene Physik. $m_e$ ist eine Ruhemasse aus dem Higgs-Yukawa-Sektor;
$m_p$ besteht zu $99\,\%$ aus QCD-Konfinementenergie. Ein Verhältnis
dieser beiden Größen verbindet zwei unabhängige Skalen
($\Lambda_{\mathrm{QCD}}$ und $y_e \cdot v$). Es gibt keinen
strukturellen Grund, warum die Geometrie eines einzelnen Parameters
$\xi$ beide Skalen gleichzeitig fixieren sollte.

\textbf{Konsequenz:} $m_p/m_e$ ist kein Ableitungsziel des aktuellen
FFGFT-Korpus. Eine scheinbar exakte Fünf-Stellen-Übereinstimmung wäre
ein Warnsignal, kein Erfolgskriterium -- denn die QCD-Beiträge tragen
prozentuale Theorieunsicherheiten.

\section{Neutrinos: der andere Grenzfall}

Im Gegensatz zum Proton sind Neutrinos ein anderer Grenzfall:
extrem kleine Ruhemassen, die vollständig im Higgs-Yukawa-Sektor liegen.

\textbf{[S]} In der FFGFT werden Neutrinos als ``fast masselose
Photonen'' behandelt (Dok.~007): Sie tragen dieselbe $(n_\theta, n_\phi)$
Wicklungsstruktur wie die geladenen Leptonen, aber mit doppelter
$\xi$-Unterdrückung. Das ergibt Ruhemassen im Sub-eV-Bereich,
konsistent mit den kosmologischen Grenzen ($\sum m_\nu < 0{,}12$~eV,
Planck 2018).

\begin{table}[h]
\centering
\caption{Neutrino-Quantenzahlen und Massenunterdrückung (Dok.~007)}
\label{tab:neutrinos}
\begin{tabular}{lccccc}
\toprule
Neutrino & $n_\theta$ & $n_\phi$ & $p_i$ & Faktor & Status \\
\midrule
$\nu_e$    & 3 & 2 & $5/2$ & $\xi^{5/2}$ & \textbf{[S]} \\
$\nu_\mu$  & 5 & 4 & $3$   & $\xi^{3}$   & \textbf{[S]} \\
$\nu_\tau$ & 9 & 5 & $25/9$& $\xi^{25/9}$& \textbf{[S]} \\
\bottomrule
\end{tabular}
\end{table}

Neutrinos belegen: der Higgs-Yukawa-Sektor hat unbegrenzte
Unterdrückung durch Potenzen von $\xi$. Das Proton liegt außerhalb
dieses Sektors.


\section{Eingerollt und ausgerollt: der strukturelle Grund der Zwei-Sektor-Trennung}

Dok.~311 unterscheidet drei Wege, wie die vierte Raumrichtung des
4D-Torus $T^4$ zur dreidimensionalen Welt wird:

\begin{center}
\begin{tabular}{lll}
\toprule
Weg & Mechanismus & Problem \\
\midrule
Kompaktifizierung (Kaluza-Klein) & vierte Richtung wird Kreis, Wicklungen = Moden
  & Radius $R$ ist freier Modul \\
Projektion (cut-and-project) & 4D-Gitter $\to$ 3D-Schnitt
  & kostet Gitterkennzahlen \\
Vierte ist kein Raum & trägt Masse, Zeit, inneren Index
  & ohne Kompaktheit kein Modenturm \\
\bottomrule
\end{tabular}
\end{center}

\textbf{[K]} FFGFT kombiniert Weg~1 und Weg~3 und bindet sie über
$\tilde{T} \cdot m = 1$: Die vierte Richtung ist \emph{eingerollt}
(kompakt) \emph{und} \emph{gebunden} (sie trägt Masse). Der Radius $R$
ist dadurch nicht frei -- er ist die Masse. Das ist der einzige Fall,
in dem eine Skala entsteht, ohne dass ein Parameter nachkommt.

\subsection{Eingerollt: Ruhemassen aus dem Yukawa-Sektor}

\textbf{[K]} Quark- und Lepton-Ruhemassen entstehen aus
\emph{eingerollten} Moden: Wicklungszahlen $(n_\theta, n_\phi)$ auf dem
kompakten Torus erzeugen quantisierte Massen
$m_i = r_i \cdot \xi^{p_i} \cdot v$ (Dok.~006, 317).
Ganzzahlige Wicklungszahlen sind erzwungen, weil der Torus kompakt ist.
Das ist der geometrische Ursprung der Yukawa-Kopplungen: keine freien
Parameter, sondern topologische Invarianten.

\subsection{Ausgerollt: Gluonen als dynamische Felder im 3D-Raum}

\textbf{[S]} Gluonen sind masselose Vektorbosonen, die im
ausgerollten, nicht-kompakten 3D-Raum propagieren. Sie sind keine
Wicklungsmoden auf dem Torus, sondern dynamische Freiheitsgrade des
Eichfeldes im entfalteten Raum. Ihre Energie kommt nicht aus
$\xi$-Wicklungen, sondern aus Bewegung und QCD-Konfinement.

\textbf{Das Proton in dieser Sprache:}
\begin{center}
\begin{tabular}{lll}
\toprule
Beitrag & Herkunft & Sektor \\
\midrule
Quark-Ruhemassen ($\sim 1\,\%$) & eingerollte Moden, Yukawa
  & FFGFT, \textbf{[K]} \\
Gluonen-Bewegungsenergie ($\sim 99\,\%$) & ausgerollte Moden, QCD-Dynamik
  & offene Brücke, \textbf{[S]} \\
\bottomrule
\end{tabular}
\end{center}

Hadronmassen haben Beiträge aus \emph{beiden} Regimen -- das ist der
strukturelle Grund, warum eine einheitliche $\xi$-Ableitung nicht
direkt greift. Die Grenze zwischen eingerollt und ausgerollt ist die
Grenze des Geltungsbereichs.

\section{Die fraktale Korrektur $K_{\text{frak}}$: Interface-Faktor beim Übergang eingerollt $\to$ ausgerollt}

\textbf{[K]} $K_{\text{frak}}$ ist der Faktor, der beim Übergang von
der geometrisch-eingerollten (bare) Beschreibung zur gemessenen
(ausgerollten, SI-) Beschreibung auftritt (Dok.~012, 133):
\begin{equation}
  m_i^{\text{SI}} = K_{\text{frak}} \cdot C_{\text{conv}} \cdot m_i^{\text{bare}},
  \qquad K_{\text{frak}} = 1 - 100\xi \approx 0{,}9867
\end{equation}
Die physikalische Begründung ist die fraktale Porosität der Raumzeit
auf Planck-Skalen ($D_{\text{frak}} = 2{,}94$, Dok.~003), die die
effektive Wechselwirkungsstärke gegenüber einem glatten Raum reduziert.

\subsection{Das Entscheidende: $K_{\text{frak}}$ kürzt sich in Verhältnissen heraus}

\textbf{[K]} In Massenverhältnissen verschwindet $K_{\text{frak}}$
vollständig (Dok.~012):
\begin{equation}
  \frac{m_\mu}{m_e} = \frac{K_{\text{frak}} \cdot m_\mu^{\text{bare}}}
  {K_{\text{frak}} \cdot m_e^{\text{bare}}} = \frac{m_\mu^{\text{bare}}}{m_e^{\text{bare}}}
\end{equation}

Das ist strukturell bedeutsam: Die Massenverhältnisse, die FFGFT im
Leptonsektor ableitet ($m_\mu/m_e$, $m_\tau/m_\mu$, Dok.~317), sind
rein geometrisch und brauchen $K_{\text{frak}}$ \emph{nicht}.
$K_{\text{frak}}$ tritt nur auf, wenn absolute SI-Werte gefragt sind
-- bei der Gravitationskonstante $G$, der Feinstrukturkonstante
$\alpha$ und absoluten Massenwerten in MeV oder kg.

\subsection{Geltungsbereich und Grenze für Hadronen}

$K_{\text{frak}}$ ist der Interface-Faktor zwischen zwei Ebenen:
\begin{center}
\begin{tabular}{lll}
\toprule
Ebene & Beschreibung & $K_{\text{frak}}$ \\
\midrule
Eingerollt (bare) & $m^{\text{bare}}$, Torus-Geometrie, $\xi$-Potenzen
  & nicht nötig \\
Übergang & $m^{\text{SI}} = K_{\text{frak}} \cdot C_{\text{conv}} \cdot m^{\text{bare}}$
  & \textbf{hier} \\
Ausgerollt (SI) & messbare Masse in MeV, kg & bereits enthalten \\
\bottomrule
\end{tabular}
\end{center}

Für den \textbf{Hadronsektor} tritt $K_{\text{frak}}$ beim Übergang
$\Lambda_{\text{QCD}}^{\text{bare}} \to \Lambda_{\text{QCD}}^{\text{SI}}$
auf -- aber $\Lambda_{\text{QCD}}^{\text{bare}}$ selbst ist nicht aus
$\xi$ abgeleitet. $K_{\text{frak}}$ als Interface-Faktor löst das
inhaltliche Problem nicht: Der Beitrag von $928$~MeV fehlt nicht wegen
fehlender Umrechnung, sondern weil der zugrunde liegende
QCD-Mechanismus im FFGFT-Korpus noch nicht behandelt ist.

\subsection{A135: Settability als Diagnose}

\textbf{[K]} A135 formuliert das Settability-Kriterium: Was sich auf
Eins setnormieren lässt, war eine Eigenschaft des Maßsystems.
$K_{\text{frak}} = 1 - 100\xi$ kann nicht auf~$1$ gesetzt werden,
ohne $\xi$ zu ändern -- es ist an etwas Physikalisches gebunden,
kein reines Maßsystem-Artefakt. Für Präzisionsvorhersagen gilt
Dok.~133.

\section{Geltungsbereich der $\xi$-Masseableitungen}

\begin{table}[h]
\centering
\caption{Vollständiger Geltungsbereich der Masseableitungen in FFGFT}
\label{tab:geltung}
\begin{tabular}{llll}
\toprule
Teilchen & Dominanter Beitrag & Status & Dok. \\
\midrule
$e, \mu, \tau$ & Higgs-Yukawa, Ruhemasse
  & \textbf{[K]} & 006, 317 \\
$\nu_e, \nu_\mu, \nu_\tau$ & Higgs-Yukawa, $\xi^2$-supprimiert
  & \textbf{[S]} & 007 \\
$u, d, s, c, b, t$ (Ruhe) & Higgs-Yukawa
  & \textbf{[S]} offen & 006 \\
$W, Z$ & Higgs-Mechanismus
  & \textbf{[S]} Skala korrekt & --- \\
Higgs $h$ & Higgs-Selbstkopplung
  & \textbf{[K]} $\xi_{\mathrm{Higgs}}$ & A-Serie \\
$p, n$ (gesamt) & QCD + Higgs-Yukawa
  & \textbf{offene Brücke} & --- \\
$m_p/m_e$ & sektorgemischt
  & \textbf{nicht beansprucht} & --- \\
\bottomrule
\end{tabular}
\end{table}

\section{Drei Szenarien für die QCD-Brücke}

Ob $\xi$ auch den QCD-Sektor erfasst, ist eine offene Forschungsfrage.
Drei Szenarien sind denkbar:

\begin{enumerate}
  \item \textbf{$\xi$ fixiert $\Lambda_{\mathrm{QCD}}$ über
    Orbifold-Geometrie.} Die $T^4/\mathbb{Z}_3$-Struktur erzeugt neun
    Orbifold-Fixpunkte. Es ist offen, ob diese eine natürliche QCD-Skala
    erzwingen. Ein Vorgriff findet sich in Dok.~041:
    $\Lambda_{\mathrm{QCD}} = v \cdot \xi^{1/3} \approx 200$~MeV
    (experimentell: $\sim 217$~MeV). Das ist ein grober Treffer,
    aber keine Herleitung -- die Formel ist dort aufgelistet,
    nicht begründet. \textbf{[S]}

  \item \textbf{Spuranomalie auf kompakten Räumen.} Die Renormierungsgruppe
    der starken Kopplung $\alpha_s(\mu)$ auf kompakten Tori erzeugt
    eine natürliche Infrarot-Skala. Khanna et al.\ (2014, arXiv:1409.1245)
    zeigen den mathematischen Rahmen; ob $\xi$ dort eingreift,
    ist nicht untersucht. \textbf{[S]}

  \item \textbf{Hadron-Sektor mit SM-Inputs (Dok.~005) und
    geometrischer $\alpha_s$-Ableitung (Dok.~160).}
    Dok.~005 enthält eine vollständige Proton-Berechnung
    ($0{,}938100$~GeV, Abweichung $0{,}02\,\%$) mit der erweiterten
    fraktalen Formel. Diese verwendet $\alpha_s = 0{,}118$
    und $\Lambda_{\mathrm{QCD}} = 0{,}217$~GeV als
    Standardmodell-Inputs. Die Aussage ``ohne freie Parameter''
    meint: keine zusätzlichen Parameter über die SM-Inputs hinaus.

    Entscheidende Präzisierung:
    Dok.~160 (Lepton-Lebensdauern) leitet $\alpha_s$ geometrisch ab:
    \begin{equation}
      \alpha_s(m_\tau) = D \cdot \xi^{1/(D+1)} = 3 \cdot \xi^{1/4}
      \approx 0{,}322 \quad (\text{exp: }0{,}330,\; -2{,}3\,\%)
    \end{equation}
    $D = 3$ Farbfreiheitsgrade, $D+1 = 4$ Raumzeit-Dimensionen.
    Das ist $\alpha_s$ bei der Skala $m_\tau$, nicht bei $M_Z$.
    Die starke Kopplung läuft mit der Energie: $\alpha_s(m_\tau)
    \approx 0{,}33$ und $\alpha_s(M_Z) \approx 0{,}118$ sind
    verschiedene Werte \emph{derselben} laufenden Kopplung.
    Die FFGFT-Ableitung trifft die richtige Skala (Dok.~160)
    und verwendet $0{,}118$ für die Proton-Rechnung in Dok.~005
    weil dort die Skala $M_Z$ relevant ist. \textbf{[K]}

    Offener Punkt: Die vollständige RGE-Brücke von
    $\alpha_s(m_\tau)$ zu $\alpha_s(M_Z)$ innerhalb der FFGFT
    ist noch nicht ausgeführt. \textbf{[S]}
\end{enumerate}

Szenario~3 ist der aktuelle Stand. Szenarien~1 und~2 sind markierte
Forschungsfragen. Der Ansatz in Dok.~041 ($\Lambda_{\mathrm{QCD}} = v
\cdot \xi^{1/3}$) ist ein erster Schritt zu Szenario~1, aber noch
keine vollständige Herleitung.

\section{Bezug zum Korpus}

\subsection{Strukturelle Einordnung der frühen Dokumente}

\textbf{[K]} Die frühen Dokumente (Dok.~001--041) haben primär die
\emph{strukturellen Zusammenhänge} zwischen $\xi$ und den physikalischen
Konstanten erarbeitet -- keine Präzisionsableitungen. Die Formeln sind
Ordnungsgrößen-Ansätze: Sie zeigen, \emph{dass} und \emph{wie} eine
Kopplungskonstante von $\xi$ abhängt, nicht notwendig den genauen
Vorfaktor oder die Einheitenumrechnung.

\begin{center}
\begin{tabular}{lll}
\toprule
Dokument & Inhalt & Charakter \\
\midrule
Dok.~041 & $\alpha_s = \xi^{-1/3}$, $\alpha_W = \xi^{1/2}$ (nat.~Einh.)
  & Strukturskizze \textbf{[S]} \\
Dok.~041 & $\Lambda_{\mathrm{QCD}} = v \cdot \xi^{1/3} \approx 200$~MeV
  & Ordnungsgröße \textbf{[S]} \\
Dok.~160 & $\alpha_s(m_\tau) = 3\xi^{1/4} \approx 0{,}322$
  & Präzisionsableitung \textbf{[K]} \\
Dok.~011 & $\alpha$ aus Torus-Geometrie
  & Präzisionsableitung \textbf{[K]} \\
Dok.~006 & Leptonmassen $m_i = r_i\xi^{p_i}v$
  & Präzisionsableitung \textbf{[K]} \\
\bottomrule
\end{tabular}
\end{center}

Dok.~041 listet $\alpha_s \approx 0{,}118$ (Experiment, bei $M_Z$) und
$\xi^{-1/3} = 9{,}65$ (natürliche Einheiten) nebeneinander -- ohne
Einheitenbrücke. Das ist kein Fehler in der Physik, sondern eine noch
nicht ausgeführte Umrechnung. Dok.~160 liefert die sauberere Ableitung
bei der Skala $m_\tau$.

\subsection{Einzelverweise}

\begin{itemize}
  \item \textbf{Dok.~005}: Proton mit erweiterter fraktaler Formel;
    $\alpha_s = 0{,}118$ und $\Lambda_{\mathrm{QCD}} = 0{,}217$~GeV
    als SM-Inputs.
  \item \textbf{Dok.~006}: Lepton-Ruhemassen [K]; Hadronmassen nicht.
  \item \textbf{Dok.~007}: Doppelte $\xi$-Unterdrückung, Neutrinos.
  \item \textbf{Dok.~041}: Strukturelle Zusammenhänge; Ordnungsgrößen,
    keine SI-Präzision.
  \item \textbf{Dok.~160}: $\alpha_s(m_\tau) = 3\xi^{1/4}$,
    Präzision $-2{,}3\,\%$ \textbf{[K]}.
  \item \textbf{Dok.~317}: Vollständige Quantenzahltabelle.
  \item \textbf{A-Serie}: $\xi_{\mathrm{Higgs}}$, Koide-Relation.
\end{itemize}

\section{Eintrag in Dok.~190 (R76)}

Diese Klarstellung wird als Registereintrag R76 in Dok.~190 deklariert:
Geltungsbereich der Masseableitungen explizit begrenzt; Hadron-Sektor
als offene Brücke dokumentiert; $m_p/m_e$ nicht als Ableitungsziel des
aktuellen Korpus beansprucht; drei Szenarien für die QCD-Brücke als
Forschungsfragen markiert.
